\chapter{Libraries of Other Languages}\label{ch-other-languages}

A package of Zeno is written in Zeno. Most of the libraries in the
world are not: Python has one for nearly everything, npm has more
packages than any other place, and Java has a very large library of
its own and a jar for whatever else. Since 0.5.2 a Zeno program can
use all three as they are, with one line at the top:

\begin{zeno}[file=mean.zeno]
import python

let statistics = python.module("statistics")
let readings = [21.5, 22.0, 23.5, 22.5]
print(statistics.mean(readings), statistics.median(readings))
\end{zeno}

\verb|python.module("statistics")| is what \verb|import statistics|
is in Python. What it gives is the module, and from there the program
is written as it would be in Python: a dot for what the module has,
brackets for a call.

The three work the same way, and this chapter takes Python first,
then says what is different about JavaScript and Java.

\begin{center}
\begin{tabular}{lP{7.4cm}}
\toprule
Written & Gives \\
\midrule
\verb|import python| & \verb|python.module(name)| and \verb|python.copy(value)| \\
\verb|import node| & \verb|node.module(name)|, \verb|node.global(name)|,
  \verb|node.new(Class, ...)| and \verb|node.copy(value)| \\
\verb|import java| & \verb|java.type(name)|, \verb|java.jar(path)| and
  \verb|java.copy(value)| \\
\bottomrule
\end{tabular}
\end{center}

\begin{note}
Nothing changes for a program that does not import one of these. They
are added beside the language, and they are there only when they are
asked for. The language that runs the library is started the first
time the program calls into it, once, and stays until the program
ends. A program that imports \verb|python| and never calls it starts
nothing.
\end{note}

\section{What it needs}

The other language must be on the computer, with the library
installed for it:

\begin{center}
\begin{tabular}{llP{6.2cm}}
\toprule
Import & Needs & Zeno looks for \\
\midrule
\verb|python| & Python 3 & \verb|python3|, then \verb|python|
  (on Windows \verb|python|, \verb|python3|, then \verb|py|);
  \verb|ZENO_PYTHON| names another \\
\verb|node| & Node 22 or later & \verb|node|; \verb|ZENO_NODE| names another \\
\verb|java| & Java 11 or later & \verb|java|, then in \verb|JAVA_HOME|;
  \verb|ZENO_JAVA| names another \\
\bottomrule
\end{tabular}
\end{center}

The Python that a virtual environment puts first, when it is active,
is the one Zeno finds. A library is looked for where that language
would look for it if the program were written in it and lay in the
same place: so a \verb|helper.py| beside the program is
\verb|python.module("helper")|, and an npm package is installed with
\verb|npm install| in the folder of the program. The program can then
be run from anywhere.

\section{Values go across in two ways}

\textbf{\texttt{none}, a bool, a number and a str go as themselves.}
What the other language returns of these kinds is a value of Zeno at
once, and what Zeno passes arrives as that language's own:

\begin{zeno}[file=day.zeno, dir=python]
import python

let datetime = python.module("datetime")
let day = datetime.date(2026, 11, 18)
print(day, day.year + 1, day.strftime("%A"))
\end{zeno}

\verb|day.year| is an int of Zeno, so \verb|+ 1| is Zeno's.
\verb|day.strftime("%A")| is a str of Zeno. Python's \verb|None| is
\verb|none|.

\textbf{Everything else stays where it was made.} \verb|datetime| and
\verb|day| above are not copies of anything. Zeno holds a number that
Python knows the value by, and each time the program writes a field, a
call or an item of it, Zeno asks Python. When nothing in the program
holds the value any more, Python is told, and lets it go.

A value that stays in the other language can be used in these ways:

\begin{center}
\begin{tabular}{lP{6.8cm}}
\toprule
Written & What the other language does \\
\midrule
\verb|thing.name| & \verb|thing.name| \\
\verb|thing.name = value| & the same \\
\verb|thing(1, 2, key=3)| & the call; what a named argument becomes is
  in the sections below \\
\verb|thing[index]| & the item, and \verb|thing[index] = value| sets it \\
\verb|len(thing)| & how many items it has \\
\verb|for item in thing { }| & every item, taken when the loop begins \\
\verb|item in thing| & whether it holds the item \\
\verb|thing == other| & that language's own test for equal \\
\verb|print(thing)|, \verb|str(thing)| & the value as that language writes it \\
\bottomrule
\end{tabular}
\end{center}

What it cannot be is the test of an \verb|if|, because a condition in
Zeno is \verb|true| or \verb|false| and nothing else. Write
\verb|len(thing) > 0| or \verb|thing is not none|. Nor can it be one
side of \verb|+| or \verb|<|: take a number of Zeno out of it first.

\begin{zeno}[file=later.zeno, dir=python]
import python

let day = python.module("datetime").date(2026, 11, 18)
let later = day.replace(month=12)
print(later.isoformat(), later == day, [day, later])
\end{zeno}

\verb|month=12| is a keyword argument of Python. Inside a list, a value
of Python is shown as Python writes it there, as a str of Zeno is shown
with its quotes inside a list.

\section{Lists and dicts}

A list of Python stays in Python too, and that is on purpose:

\begin{zeno}[file=names.zeno, dir=python]
import python

let builtins = python.module("builtins")
let names = builtins.list(["cy", "ann"])
names.append("bob")
names.sort()
print(names, len(names), names[0], "bob" in names)
for name in names {
    print(name.upper())
}
\end{zeno}

\verb|names| is the one list there is. \verb|append| and \verb|sort|
are Python's, and they change the list that Python has. If the list
came across as a copy, \verb|names.append("bob")| would add to the
copy, Python would never see it, and nothing would say so. Each
\verb|name| in the loop is a str of Zeno, so \verb|upper| is Zeno's.

When a list of Zeno is what the program wants, to pass to a function
that takes \verb|list[str]| or to use Zeno's own methods on, the
\verb|copy| of the module makes one:

\begin{zeno}[file=mine.zeno, dir=python]
import python

let names = python.module("builtins").sorted(["cy", "ann", "bob"])
let mine = python.copy(names)
mine.append("dee")
print(mine, len(names))
\end{zeno}

\verb|copy| takes a list, a tuple or a dict of Python, and gives a list
or a dict of Zeno, with every list and dict inside it copied as well.
A dict whose keys are not all str or int cannot be a dict of Zeno, and
\verb|copy| raises a \verb|ValueError| that names the key.

\textbf{The other way is always a copy.} A list or dict of Zeno passed
to Python arrives as a new list or dict there:

\begin{zeno}[file=shuffle.zeno, dir=python]
import python

let xs = [3, 1, 2]
python.module("random").shuffle(xs)
print(xs)        # Python shuffled its copy
\end{zeno}

To have Python work on a list and see the result, make the list in
Python (\verb|builtins.list(xs)|) and copy it back.

\section{Errors}

An error that the other language raises is an error of Zeno, and
\verb|try| catches it (\hyperref[ch-errors]{Errors}):

\begin{zeno}[file=oops.zeno, dir=python]
import python

let json = python.module("json")
try {
    json.loads("{oops")
} except ValueError as error {
    print(error.message)
}
\end{zeno}

The message begins with what Python called the error when that is not
the name of the kind. The kinds are matched like this:

\begin{center}
\begin{tabular}{lP{8.4cm}}
\toprule
From & Becomes \\
\midrule
Python & \verb|OSError| and its kinds are a \verb|FileError|;
  \verb|ValueError|, \verb|IndexError|, \verb|KeyError|,
  \verb|OverflowError| and \texttt{Zero\-Division\-Error} are the kind
  of the same name \\
JavaScript & an error the system gave (no such file, no network) is a
  \verb|FileError| \\
Java & \verb|IOException| and its kinds are a \verb|FileError|;
  \texttt{Index\-OutOf\-Bounds\-Exception} is an \verb|IndexError|;
  \texttt{Illegal\-Argument\-Exception} is a \verb|ValueError| \\
\bottomrule
\end{tabular}
\end{center}

\noindent Anything else is an \verb|Error|. Two things can go wrong
before any library is reached, and both are a \verb|FileError|: there
is no Python (or Node, or Java) to run, or it has no module of that
name. The message says which one Zeno found, and where it looked.

\begin{zeno}[file=missing.zeno, dir=python, fails=yes]
import python

let video = python.module("no_such_library")
\end{zeno}

What another language can raise is not known before the program runs.
A function that declares what it raises (\hyperref[ch-errors]{Errors})
and calls into one says \verb|raises Error|, which is every kind.

\section{Types}

What a module has is known when the program runs, and not before. So
a value from another language has the type that \verb|json.parse|
gives (\hyperref[ch-reading-writing]{Reading and Writing}): the checker
lets everything through, and a wrong name is an error at its line when
the program gets there. Give a value a type of Zeno where it lands,
and from there on it is checked:

\begin{zenofrag}
let title: str = about["title"]       # checked when it happens
\end{zenofrag}

For the same reason \verb|zeno check --strict| refuses a program that
imports one of the three.

\section{JavaScript}

\verb|node.module(name)| is what \verb|require(name)| is in
JavaScript. A module that comes with Node, such as \verb|fs|,
\verb|path|, \verb|os|, \verb|crypto| or \verb|zlib|, needs nothing
more; a package from npm is installed in the program's folder with
\verb|npm install|. A file of JavaScript beside the program is a
module by its path, \verb|node.module("./helper.js")|.

Much of JavaScript is not in a module at all. \verb|JSON|,
\verb|Math|, \verb|Date| and \verb|process| are there everywhere, and
\verb|node.global| gives them. JavaScript makes an object of a class
with \verb|new|, and Zeno with \verb|node.new|:

\begin{zeno}[file=globals.zeno]
import node

let path = node.module("path")
print(path.basename("/home/ann/notes.txt"), path.extname("notes.txt"))

let JSON = node.global("JSON")
let about = JSON.parse("{\"name\": \"ann\", \"age\": 31}")
print(about.name, about.age + 1, about.email is none)

let start = node.new(node.global("Date"), 2026, 10, 18)
print(start.getFullYear(), start.toDateString())
\end{zeno}

\verb|about.email| is \verb|none|, and not an error: a field that an
object of JavaScript does not have is \verb|undefined|, and both
\verb|null| and \verb|undefined| are \verb|none| in Zeno.

Four things are different from Python.

\begin{itemize}
\item \textbf{One kind of number.} A number of JavaScript that is
  whole arrives as an int, and any other as a float. So
  \verb|Math.sqrt(16.0)| is \verb|4|, an int; where a float is
  needed, \verb|float(x)| makes one.
\item \textbf{Named arguments are the options.} JavaScript has no
  named arguments. A function there takes its options as one object,
  its last argument, and so the named arguments of a call are gathered
  into that object: \verb|fs.readFileSync("notes.txt", encoding="utf8")|
  is \verb|readFileSync("notes.txt", {encoding: "utf8"})|.
\item \textbf{What takes time is waited for.} A call that gives a
  promise gives its answer, and the program goes on when it is there.
  There is no \verb|await| to write, and a promise that fails is an
  error at the line of the call.
\item \textbf{Equal is the same object.} \verb|==| is JavaScript's
  \verb|===|, so two arrays with the same items are not equal. Compare
  their copies: \verb|node.copy(a) == node.copy(b)|.
\end{itemize}

\begin{zeno}[file=waited.zeno]
import node

let Math = node.global("Math")
print(Math.max(3, 9, 4), Math.sqrt(16.0), Math.sqrt(2.0))

let options = node.global("Object")(recursive=true, retries=3)
print(options, options.retries + 1)

let crypto = node.module("crypto")
print(crypto.createHash("sha256").update("zeno").digest("hex"))
\end{zeno}

\verb|node.global("Object")| called with named arguments makes such an
object by itself, to keep or to add to. \verb|node.copy| takes an
array or a \verb|Set| and gives a list, and a plain object or a
\verb|Map| and gives a dict.

\section{Java}

\verb|java.type(name)| is what \verb|import| of a class is in Java,
and the name is always the whole name that Java's own documentation
gives. A dot after a class is what is \verb|static| in it, and calling
a class makes an object, as \verb|new| does there:

\begin{zeno}[file=classes.zeno]
import java

let Math = java.type("java.lang.Math")
print(Math.max(3, 9), Math.floorMod(-7, 3))

let LocalDate = java.type("java.time.LocalDate")
let day = LocalDate.of(2026, 11, 18)
print(day, day.getDayOfWeek(), day.plusMonths(1), day.getYear() + 1)

let names = java.type("java.util.ArrayList")(["cy", "ann"])
names.add("bob")
java.type("java.util.Collections").sort(names)
print(names, len(names), names[0], "bob" in names, names == ["ann", "bob", "cy"])
\end{zeno}

A list of Zeno passed to a class of Java arrives as a new
\verb|ArrayList|, and a dict as a \verb|LinkedHashMap|: so
\verb|ArrayList(["cy", "ann"])| is the short way to make a collection
of Java from what Zeno has. \verb|==| is Java's \verb|equals|, so a
list of Java equals a list of Zeno with the same items.

What is different from Python:

\begin{itemize}
\item \textbf{A library is a jar.} \verb|java.jar("lib/gson-2.11.0.jar")|
  says where one is, counted from the folder of the \verb|.zeno| file,
  and after it the classes of that jar are found. A folder in place of
  a jar is every jar in it. Kotlin, Scala and Clojure libraries are
  jars as well.
\item \textbf{Java's method decides.} Java can have several methods of
  one name that take different types. The one that the arguments fit
  best is called, in the order Java itself would choose: an int is an
  \verb|int| before it is a \verb|long|, and a \verb|long| before a
  \verb|double|. When none fits, the error lists what there is.
\item \textbf{No named arguments.} Java has none, so a call with one is
  an error that says so.
\item \textbf{The first run takes a second longer.} Zeno's own part on
  the Java side is compiled once, with \verb|javac|, and kept in the
  folder for what can be made again (\verb|~/.cache/zeno| on Linux).
\end{itemize}

\begin{zeno}[file=valueof.zeno]
import java

let String = java.type("java.lang.String")
print(String.valueOf(5), String.valueOf(2.5), String.valueOf(true))

try {
    java.type("java.lang.Integer").parseInt("forty")
} except ValueError as error {
    print(error.message)
}
\end{zeno}

Those are three different methods of Java, and Zeno picked each from
the type of what it was given.

\begin{note}
One trap is Java's own. A \verb|List| has \verb|remove(int)|, which
takes out the item at a place, and \verb|remove(Object)|, which takes
out an item. \verb|list.remove(30)| is the first, in Java as here. To
take the number 30 out of a list of numbers, write
\verb|list.remove(list.indexOf(30))|.
\end{note}

\section{More than one in a program}

A program can import all three. Each value stays where it was made. A
str or a number that one gave can be passed to another, because by
then it is a value of Zeno; a value that stayed in one cannot be given
to another, and Zeno says so at the line:

\begin{zeno}[file=both.zeno, fails=yes]
import python
import node

let day = python.module("datetime").date(2026, 11, 18)
let JSON = node.global("JSON")
print(JSON.stringify(day.isoformat()))
print(JSON.stringify(day))
\end{zeno}

The first \verb|print| works, because \verb|isoformat()| gave a str of
Zeno. The \verb|copy| of each module makes a list or a dict of Zeno,
and that can go anywhere.

\section{A name that is a word of Zeno}

Python has \verb|re.match|, JavaScript has \verb|Array.from|, and Java
has \verb|System.in|; \verb|match|, \verb|from| and \verb|in| are
words of Zeno, so none of them parses after a dot. Each language has
its own way around it:

\begin{center}
\begin{tabular}{lP{8.4cm}}
\toprule
Language & Write \\
\midrule
Python & \verb|builtins.getattr(re, "match")(p, text)| \\
JavaScript & \verb|Array["from"](items)|: brackets are a dot there \\
Java & \verb|System["in"]|, \verb|LocalDate["from"](other)| \\
\bottomrule
\end{tabular}
\end{center}

\section{What is not there yet}

\begin{itemize}
\item \textbf{A function of Zeno cannot be given to another language.}
  A library that calls back, for progress, a sort key, an event or a
  request to a server, cannot call into Zeno yet, and Zeno says so at
  the line. Much of Node gives a promise instead, and that works.
\item \textbf{A loop takes everything first.} \verb|for line in thing|
  asks for all the items, then loops. Over something that never ends,
  it never begins.
\item \textbf{A class of Zeno cannot extend a class of Java}, or be
  one of its interfaces.
\item \textbf{Built programs.} \verb|zeno run| and \verb|zeno test| have
  all three, with and without \verb|--vm|. \verb|zeno build|, in every
  form and for every target, refuses them, because a built program
  must run where the other language may not be.
\end{itemize}

\section{How it is built}

The other language runs as a program of its own, beside Zeno. A small
helper for it is part of the \verb|zeno| file, and the first call
starts the language with it. Zeno opens a socket that only this
computer can reach, and gives the helper its port and a word made up
for this run; the helper connects, says the word, and only then is it
listened to. Because the conversation has a line of its own, the
other language keeps the screen and the keyboard: its \verb|print|, a
progress bar that redraws itself and \verb|input()| all work as if it
had been started by hand.

Each question is one line of JSON, and so is each answer:

\begin{zenofrag}[lang=shell]
{"op":"get","on":["p",3],"name":"year"}                 day.year
{"ok": 2026}
{"op":"call","on":["p",9],"args":["{oops"],"named":[],"free":[5,7]}
{"err": ["ValueError", "JSONDecodeError: Expecting ..."]}
\end{zenofrag}

\verb|["p", 3]| is the value the helper keeps under the number 3.
\verb|"free"| rides along with the next question: the numbers Zeno has
let go of since the last one. In Zeno the value is one more kind of
\verb|Value|, which holds the number, and the operations that both
engines share (a field, an item, \verb|len|, \verb|==|) ask over the
socket when they meet it. Nothing on Zeno's side knows Python but the
name of the helper, so JavaScript and Java are each another helper that
speaks the same lines. The reasons are chapters 26, 28 and 29 of
\verb|docs/decisions.md|.

\section{Things to try}

\begin{enumerate}
\item Run the first program of this chapter. Then misspell
  \verb|statistics| and read the message. Which Python does it name?
\item Print \verb|python.module("sys").executable|. Start a virtual
  environment and run it again.
\item \verb|builtins.sorted([3, 1, 2])| gives a list that stays in
  Python. Print it, then print \verb|python.copy| of it. What is
  different about the two lines, and why?
\item Write \verb|function mean(xs: list[float]) -> float| in Zeno, and pass
  it a list of Python with and without \verb|copy|. What do the
  checker and the program say?
\item With Node, read a file through \verb|node.module("fs/promises")|.
  Where is the promise?
\item With Java, call \verb|java.type("java.util.ArrayList")().get(3)|
  inside a \verb|try|. Which kind of error catches it?
\end{enumerate}
